\chapter{Prediction and Betting Markets}\label{ch:m3:prediction-and-betting-markets}

In September 2024 a US court ruled that the derivatives regulator had exceeded its authority in
forbidding a regulated exchange to list contracts on which party would control Congress. Two years
earlier the same regulator had fined a \omterm{def:m3:blockchains-for-traders:chain}{blockchain} venue that offered, without registration, contracts
such as ``Will Trump win the 2020 presidential election?''. A contract that pays one dollar if an event happens trades at a price that reads as a
probability. \omterm{def:m3:prediction-and-betting-markets:book}{Bookmakers} have quoted such prices for centuries, with a margin; exchanges now let bettors
trade them with each other, and derivatives regulators have had to decide which of them are gambling.
This chapter explains how odds encode probabilities and margins, how \omterm{def:m3:prediction-and-betting-markets:exchange}{betting exchanges} and \omterm{def:m3:prediction-and-betting-markets:event}{event
contracts} work, how trading firms make these markets, and the best-documented bias in them.

\section{Bookmakers, overround and implied probabilities}

\begin{definition}[Bookmaker, implied probability, overround]\label{def:m3:prediction-and-betting-markets:book}
A \emph{bookmaker}\index{bookmaker} is a firm that quotes odds on the outcomes of an event and takes bets against them as the
counterparty. The \emph{implied probability}\index{implied probability} of decimal odds $o$ is $1/o$. The \emph{overround}\index{overround}
of a book is the amount by which its implied probabilities over all outcomes exceed one: the bookmaker's margin.
\end{definition}

Decimal odds $o$ pay $o$ per unit staked, the stake included; they are the odds of One Quant Book 2, chapter 29, written as a
payout. A three-way market at 1.50, 4.00 and 7.00 has implied probabilities 66.7\%, 25.0\% and 14.3\%, summing to 105.95\%: an
\omterm{def:m3:prediction-and-betting-markets:book}{overround} of 5.95\%. A bettor who backs every outcome in proportion to its \omterm{def:m3:prediction-and-betting-markets:book}{implied probability} loses 5.6\% of the stake whatever
happens. To read the market's probabilities, the margin must be taken out, and how to take it out is a model of where the \omterm{def:m3:prediction-and-betting-markets:book}{bookmaker}
put it.

\begin{proposition}[Four ways to remove the margin]\label{prop:m3:prediction-and-betting-markets:removal}
Let $\pi_i = 1/o_i$ with $\sum \pi_i = 1 + m$. The multiplicative method sets $p_i = \pi_i/(1 + m)$; the additive method $p_i = \pi_i -
m/n$; the power method $p_i = \pi_i^k$ with $k > 1$ such that $\sum p_i = 1$; Shin's method solves for the share $z$ of bettors who are
insiders, with $p_i = \big(\sqrt{z^2 + 4(1 - z)\pi_i^2/(1 + m)} - z\big)/(2(1 - z))$ and $\sum p_i = 1$. The multiplicative method puts the
margin in proportion to the probability; the others put relatively more of it on the longshots.
\end{proposition}

\begin{proof}
Each formula sums to one by construction (the power and Shin parameters by a one-dimensional root search, unique since the sums are
monotone in the parameter). For the power method, $\pi_i^k/\pi_i = \pi_i^{k - 1}$ is smallest for the smallest $\pi_i$, so longshots
lose the most probability relative to their implied value; Shin's formula does the same because the insider term penalises small
$\pi_i$ proportionally more.
\end{proof}

\begin{table}[ht]
\centering
\small
\begin{tabular}{@{}lrrr@{}}
\toprule
Method & favourite (1.50) & middle (4.00) & longshot (7.00) \\
\midrule
Implied, with margin & 66.67 & 25.00 & 14.29 \\
Multiplicative & 62.92 & 23.60 & 13.48 \\
Additive & 64.68 & 23.02 & 12.30 \\
Power & 64.82 & 22.71 & 12.48 \\
Shin ($z = 3.0\%$) & 64.23 & 23.15 & 12.62 \\
\bottomrule
\end{tabular}
\caption{Probabilities, in percent, recovered from a three-way book with an \omterm{def:m3:prediction-and-betting-markets:book}{overround} of 5.95\% by the methods of
\cref{prop:m3:prediction-and-betting-markets:removal}.}\label{tab:m3:prediction-and-betting-markets:methods}
\end{table}

\section{Betting exchanges}

\begin{definition}[Betting exchange, back bet, lay bet, in-play betting]\label{def:m3:prediction-and-betting-markets:exchange}
A \emph{betting exchange}\index{betting exchange} is a venue on which bettors trade bets with each other through an order book and the
operator charges a commission rather than taking positions. A \emph{back bet}\index{back bet} is a bet that an outcome happens; a \emph{lay
bet}\index{lay bet} is a bet that it does not, the layer taking the role of the \omterm{def:m3:prediction-and-betting-markets:book}{bookmaker} for that bet. \emph{In-play betting}\index{in-play
betting} is betting while the event is under way, at prices that move with it.
\end{definition}

A \omterm{def:m3:prediction-and-betting-markets:exchange}{betting exchange} is a limit order book (One Quant Book 1, chapter 19) in which the instrument is an outcome and the price is decimal odds.
Backers want high odds, layers low ones, and the spread between the best back and lay odds is the market's spread. Commission is charged on each
market's net winnings, so a bettor pays only on what it wins: Betfair computes it as net winnings times the market's base rate times one minus
the bettor's discount, with base rates in its Australian help pages of 6\% for sport, 10\% for one rugby league competition and 8 or 10\% for local
racing. Lay and back positions offset: a bettor who backed at 3.0 and lays at 2.5 after the
odds shorten locks in a profit whatever happens, just as a trader closes a long at a higher price. In-play markets move with every event on the
field and reward speed of information, which puts slow layers at the mercy of fast backers.

\section{Event contracts on and off chain}

\begin{definition}[Event contract, resolution source]\label{def:m3:prediction-and-betting-markets:event}
An \emph{event contract}\index{event contract} is a derivative that pays a fixed amount, usually one dollar, if a specified event occurs and nothing
otherwise; its price, as a share of the payout, reads as the market's probability of the event. Its \emph{resolution source}\index{resolution source}
is the source and procedure named in the contract by which the occurrence of the event is determined.
\end{definition}

In the United States \omterm{def:m3:prediction-and-betting-markets:event}{event contracts} are derivatives under the Commodity Exchange Act and must be listed on a \omterm{def:m3:getting-access-commodities-and-power:dcm}{designated contract market}
(\cref{ch:m3:getting-access-commodities-and-power}). The statute's special rule lets the regulator prohibit \omterm{def:m3:prediction-and-betting-markets:event}{event contracts} involving certain
activities, among them unlawful activity and gaming, if contrary to the public interest.

\begin{dated}{2026-09}{Event contracts before US courts and regulators}\label{dat:m3:prediction-and-betting-markets:kalshi}
On 12 September 2024 the US District Court for the District of Columbia, in KalshiEX LLC v.\ CFTC, found that Kalshi's congressional control contracts
(``will this chamber of Congress be controlled by this party for this term?'') involve neither unlawful activity nor gaming, that the CFTC's order
prohibiting them exceeded its statutory authority, and granted Kalshi summary judgment. The CFTC appealed; on 7 May 2025 the Court of Appeals granted
its unopposed motion to dismiss the appeal. In January 2022 the CFTC had settled charges against Blockratize, Inc., doing business as Polymarket, for
offering off-exchange event-based binary options without being a \omterm{def:m3:getting-access-commodities-and-power:dcm}{designated contract market} or a registered swap execution facility; the order imposed
a USD~1.4 million penalty and required it to wind down its non-compliant markets. On 2 September 2025 CFTC staff granted a no-action position
on swap reporting to QCX, a \omterm{def:m3:getting-access-commodities-and-power:dcm}{designated contract market} and clearing house that Polymarket had acquired for its return to the United States.
\end{dated}

On a \omterm{def:m3:blockchains-for-traders:chain}{blockchain} venue the \omterm{def:m3:prediction-and-betting-markets:event}{event contract} is a pair of \omterm{def:m3:blockchains-for-traders:key}{tokens}, YES and NO, minted together against one dollar of \omterm{def:m3:blockchains-for-traders:stable}{stablecoin} and redeemable for it,
traded on an order book or a market maker, and settled by an \omterm{def:m3:blockchains-for-traders:contract}{oracle} (\cref{ch:m3:blockchains-for-traders}). The \omterm{def:m3:prediction-and-betting-markets:event}{resolution source} is the contract's
most important and least visible term: ``who won the election'' must be decided by someone, at some time, from some source, and the rules for disputed
or ambiguous outcomes decide what the \omterm{def:m3:blockchains-for-traders:key}{tokens} pay. Minting also inflates volume: from the full on-chain ledger of Polymarket's 2024 presidential
markets, Tsang and Yang measure October's turnover in the Trump market at USD~391 million against 958 million of naive volume, which counts
the creation and destruction of pairs as trades.

\begin{dated}{2026-09}{An optimistic oracle}\label{dat:m3:prediction-and-betting-markets:uma}
UMA's documentation describes its optimistic \omterm{def:m3:blockchains-for-traders:contract}{oracle}: third-party proposers post a resolution with a bond, refundable if the proposal is correct; the
proposal is pending for a challenge period during which anyone may dispute it by posting a dispute bond; disputes go to a vote of UMA \omterm{def:m3:blockchains-for-traders:key}{token} stakers,
with a 24-hour commit period and a 24-hour reveal period, and the incorrect side's bond compensates the correct side.
\end{dated}

\begin{omfigure}
\centering
\begin{tikzpicture}[font=\small,
  s/.style={draw,thick,rounded corners=2pt,align=center,minimum height=0.9cm,text width=2.5cm},
  arr/.style={-{Stealth[length=2.2mm]},thick}]
\node[s,draw=omDef,fill=omDef!8] (u) at (0,0) {one dollar of\\stablecoin};
\node[s,draw=omThm,fill=omThm!8] (m) at (3.6,0) {mint one YES\\and one NO};
\node[s,draw=omProp,fill=omProp!8] (t) at (7.2,0) {trade on a book\\or a pool};
\node[s,draw=omExo,fill=omExo!8] (o) at (10.8,0) {oracle resolves:\\winner pays \$1};
\draw[arr] (u) -- (m); \draw[arr] (m) -- (t); \draw[arr] (t) -- (o);
\draw[arr,dashed] (m.south) -- ++(0,-0.6) -| (u.south);
\node[font=\footnotesize,anchor=north] at (1.8,-1.1) {a YES and a NO together redeem for one dollar at any time};
\end{tikzpicture}

\omcaption{An \omterm{def:m3:prediction-and-betting-markets:event}{event contract} on a \omterm{def:m3:blockchains-for-traders:chain}{blockchain} venue: YES and NO \omterm{def:m3:blockchains-for-traders:key}{tokens} are minted together against one dollar, trade separately, and the winning \omterm{def:m3:blockchains-for-traders:key}{token} pays
one dollar when the \omterm{def:m3:blockchains-for-traders:contract}{oracle} resolves the event. Because a pair always redeems for one dollar, YES and NO prices sum to about one. Schematic.}\label{fig:m3:prediction-and-betting-markets:tokens}
\end{omfigure}

\section{How trading firms make these markets}

A market maker in \omterm{def:m3:prediction-and-betting-markets:event}{event contracts} quotes a probability with a spread, like any other market maker, and its risks are the same in new clothes: adverse
selection by bettors who know more (One Quant Book 1, chapter 19), inventory that it cannot hedge except in the same event on other venues, and prices
that jump to zero or one at resolution. Three features are peculiar. The same event trades on several venues, a \omterm{def:m3:prediction-and-betting-markets:book}{bookmaker}, an exchange, a regulated
event-contract exchange and a \omterm{def:m3:blockchains-for-traders:chain}{blockchain} venue, with different fees, participants and \omterm{def:m3:prediction-and-betting-markets:event}{resolution sources}, so a firm's first edge is cross-venue:
buying the event where it is cheap and selling it where it is dear. Capital is locked until resolution, often for months, so a price gap must beat the
return on that capital. And the \omterm{def:m3:prediction-and-betting-markets:event}{resolution sources} differ: the same headline event can resolve differently on two venues if their sources or rules for
edge cases differ, which turns an arbitrage into a bet on the rules.

\begin{proposition}[When a cross-venue gap pays]\label{prop:m3:prediction-and-betting-markets:gap}
Let venue A sell YES at $y_A$ and venue B sell NO at $1 - y_B$ with $y_B > y_A$, with fees $f_A$ and $f_B$ per contract and capital cost $r$ a year for
$\tau$ years until resolution. Buying one of each pays exactly one dollar whatever happens, if both venues resolve the same way, and earns
$1 - (y_A + 1 - y_B + f_A + f_B)(1 + r\tau)$. It pays when the gap $y_B - y_A$ exceeds approximately $f_A + f_B + r\tau$.
\end{proposition}

\begin{proof}
One YES and one NO on the same event pay one dollar together; the cost is the two prices, the fees and the financing until the payment.
\end{proof}

\section{The favourite--longshot bias}

\begin{definition}[Favourite--longshot bias]\label{def:m3:prediction-and-betting-markets:flb}
The \emph{favourite--longshot bias}\index{favourite--longshot bias} is the tendency, documented across betting markets, for bets on longshots to return
less per unit staked than bets on favourites, so that longshots' prices overstate their probabilities.
\end{definition}

Snowberg and Wolfers, studying horse racing, asked whether the bias comes from bettors who love risk or from bettors who misperceive probabilities, and
found that misperceptions, as prospect theory predicts, explain it better. Shin's model, in which the \omterm{def:m3:prediction-and-betting-markets:book}{bookmaker} protects itself against insiders by shading
longshots most, is another explanation, and the reason the Shin method of \cref{prop:m3:prediction-and-betting-markets:removal} exists. For a trader the
bias is a structural edge on one side: selling longshots and buying favourites, at scale and with discipline, and a warning that a market's longshot prices
are poor probability estimates.

\begin{omfigure}
\begin{tikzpicture}
\begin{axis}[width=11cm,height=5cm,ybar,bar width=14pt,
  xlabel={decimal odds},ylabel={return per unit staked, \%},
  tick label style={font=\small},label style={font=\small},
  xmin=0.4,xmax=7.6,xtick={1,2,3,4,5,6,7},xticklabels={1--2,2--3,3--5,5--10,10--20,20--50,50--1000},
  ymin=-45,ymax=0,grid=major,grid style={gray!25},
  table/col sep=comma]
\addplot[fill=omThm!60,draw=omThm] table[x=idx,y=ret]{figdata/markets-3/27-prediction-and-betting-markets/flb.csv};
\end{axis}
\end{tikzpicture}

\omcaption{Realised return per unit staked by odds bucket on a synthetic book of 20\,000 eight-runner races whose \omterm{def:m3:prediction-and-betting-markets:book}{bookmaker} prices implied probabilities
proportional to the true probability raised to the power 0.92, with a 10\% margin. Favourites lose a few percent; longshots between 50 and 1\,000 lose over 40\%.
Synthetic results. Data: the chapter's tutorial.}\label{fig:m3:prediction-and-betting-markets:flb}
\end{omfigure}

\section{Tutorial: odds, margins and the bias}\label{tut:m3:prediction-and-betting-markets:odds}

\begin{tutorial}
\textbf{Goal.} Convert \omterm{def:m3:prediction-and-betting-markets:book}{bookmaker} odds to probabilities with the multiplicative, additive, power and Shin methods, compare them, and measure a
\omterm{def:m3:prediction-and-betting-markets:flb}{favourite--longshot bias} on a synthetic book of results. \textbf{End state:} \cref{tab:m3:prediction-and-betting-markets:methods},
\cref{fig:m3:prediction-and-betting-markets:flb} and the numbers of the weekend problem.
\begin{tutsteps}
\item \textbf{Shin's method.} A one-dimensional root search for the insider share.
\omcode{code/firm/odds/firm_odds.py}{52}{61}{Shin's margin removal.}\label{lst:m3:prediction-and-betting-markets:shin}
\item \textbf{The bias.} Price races with a biased \omterm{def:m3:prediction-and-betting-markets:book}{bookmaker}, draw the winners, and average the returns by odds.
\omcode{code/markets-3/27-prediction-and-betting-markets/python/m3_betting.py}{27}{44}{A synthetic book of races and its bets' realised returns.}\label{lst:m3:prediction-and-betting-markets:flb}
\item \textbf{Run} \texttt{methods()}, \texttt{returns\_by\_odds(flb\_book())}, \texttt{election\_gap()} and \texttt{fig\_betting.py}.
\end{tutsteps}
\textbf{What to change next.} Recover the bias parameter from the synthetic results by maximum likelihood; hedge a \omterm{def:m3:prediction-and-betting-markets:exchange}{back bet} to an equal outcome with
\texttt{hedge\_equal}; add an in-play price path and compute a market maker's inventory risk.
\end{tutorial}

\section{Build: odds and hedging}\label{bld:m3:prediction-and-betting-markets:odds}

\begin{build}
\bfield{Purpose} The miniature firm prices events across \omterm{def:m3:prediction-and-betting-markets:book}{bookmakers}, exchanges and event-contract venues: it needs their prices as probabilities without
margins, their fees on the same basis, and hedges that equalise its outcomes.
\bfield{Interface} \texttt{implied}, \texttt{overround}; \texttt{multiplicative}, \texttt{additive}, \texttt{power}, \texttt{shin}; \texttt{kelly(p, odds)};
\texttt{exchange\_payout(stake, odds, commission)}; \texttt{hedge\_equal(back\_stake, back\_odds, lay\_odds, commission)}.
\bfield{Rules} Decimal odds internally, conversions at the edges; probabilities that sum to one; commission applied on net winnings as exchanges do.
\bfield{Acceptance tests} \texttt{code/firm/odds/tests/}: the \omterm{def:m3:prediction-and-betting-markets:book}{overround}; all methods summing to one; power and Shin shading the longshot more than the
multiplicative method; Kelly at zero edge and at a stated edge; a hedge to equal outcomes.
\bfield{Stretch} Multi-venue arbitrage scanner for one event; the resolution-risk discount between venues with different sources.
\end{build}

\begin{omsources}
\item KalshiEX LLC v.\ CFTC, No.\ 23-3257 (D.D.C.), memorandum opinion of 12 September 2024 (Doc 51); D.C. Circuit No.\ 24-5205, order of 7 May 2025
granting the CFTC's unopposed motion to dismiss.
\item CFTC, press release 8478-22 (Polymarket), January 2022; CFTC Letter No.\ 25-28, 2 September 2025; CoinDesk, 3 September 2025.
\item Betfair, help pages on commission (``Exchange: What is Commission and how is it calculated?''; Australian ``Commissions and Charges''),
Internet Archive copies.
\item K.~P.~Tsang and Z.~Yang, ``The Anatomy of a Blockchain Prediction Market: Polymarket in the 2024 U.S. Presidential Election'',
arXiv:2603.03136, version of 14 September 2026.
\item UMA documentation, ``How does UMA's oracle work''. Accessed September 2026.
\item H.~S.~Shin, ``Measuring the incidence of insider trading in a market for state-contingent claims'', \emph{Economic Journal} 103, 1993; E.~Snowberg
and J.~Wolfers, ``Explaining the Favorite-Longshot Bias: Is it Risk-Love or Misperceptions?'', NBER Working Paper 15923, 2010.
\end{omsources}

\section{Exercises}

\begin{exercise}[$\star$]\label{exo:m3:prediction-and-betting-markets:1}
A two-way market is priced at 1.80 and 2.10. What is its \omterm{def:m3:prediction-and-betting-markets:book}{overround}, and what are the multiplicative probabilities?
\end{exercise}

\begin{exercise}[$\star$]\label{exo:m3:prediction-and-betting-markets:2}
A bettor backed at 3.0 for 100 and the odds have shortened to 2.5. How much should it lay to lock in an equal profit, and what is that profit?
\end{exercise}

\begin{exercise}[$\star$]\label{exo:m3:prediction-and-betting-markets:3}
Why must a US exchange listing \omterm{def:m3:prediction-and-betting-markets:event}{event contracts} be a \omterm{def:m3:getting-access-commodities-and-power:dcm}{designated contract market}?
\end{exercise}

\begin{exercise}[$\star\star$]\label{exo:m3:prediction-and-betting-markets:4}
Compute the power-method probabilities of the book 1.50, 4.00, 7.00, and explain why the longshot's is below the multiplicative one.
\end{exercise}

\begin{exercise}[$\star\star$]\label{exo:m3:prediction-and-betting-markets:5}
An \omterm{def:m3:prediction-and-betting-markets:event}{event contract} trades at 0.58 and you believe the probability is 0.65. What Kelly fraction should you stake, ignoring fees?
\end{exercise}

\begin{exercise}[$\star\star$]\label{exo:m3:prediction-and-betting-markets:6}
Why does an arbitrage between two venues' contracts on ``the same'' event carry resolution risk?
\end{exercise}

\begin{exercise}[$\star\star\star$]\label{exo:m3:prediction-and-betting-markets:7}
\emph{Coding.} With \texttt{returns\_by\_odds}, what is the return of backing every runner priced between 1 and 3, and of backing every runner priced between
50 and 1\,000?
\end{exercise}

\begin{exercise}[$\star\star\star$]\label{exo:m3:prediction-and-betting-markets:8}
\emph{Find the flaw.} ``The market prices the longshot at 50 to 1, so its probability is 2\%.''
\end{exercise}

\section{Problem: Election Night}

\begin{problem}[{Weekend problem --- one event, two venues}]\label{pb:m3:prediction-and-betting-markets:1}
A month before an election, venue A (a regulated event-contract exchange) sells YES on a candidate at 0.58 with a fee of 1 cent a contract; venue B (a
\omterm{def:m3:blockchains-for-traders:chain}{blockchain} venue) sells YES at 0.62, so NO at 0.38, with a fee of half a cent. Capital earns 4.5\% a year elsewhere. Your model gives the candidate 65\%.

\medskip\noindent\textbf{Part I --- The arbitrage.}
\begin{enumerate}
\item What does one YES on A and one NO on B cost, with fees?
\item What does the pair pay, and when?
\item What is the profit per pair after the cost of capital?
\item What is the smallest gap that survives fees and capital?
\item What does 10\,000 pairs earn, and what capital does it lock?
\end{enumerate}

\medskip\noindent\textbf{Part II --- The risks.}
\begin{enumerate}[resume]
\item What happens if the two venues resolve differently?
\item What does venue B's settlement in \omterm{def:m3:blockchains-for-traders:stable}{stablecoins} add?
\item What limits the size of the trade on each venue?
\item Can a US-resident firm trade on both venues?
\item What does the gap tell you about the two venues' participants?
\end{enumerate}

\medskip\noindent\textbf{Part III --- The view.}
\begin{enumerate}[resume]
\item With a belief of 65\%, what Kelly fraction would you stake on YES at venue A?
\item What is the expected return per contract at 0.58?
\item Why would a professional stake a fraction of Kelly?
\item How would you check whether 65\% is better than the market's 58--62\%?
\item How does the \omterm{def:m3:prediction-and-betting-markets:flb}{favourite--longshot bias} bear on the view?
\end{enumerate}

\medskip\noindent\textbf{Part IV --- Judgement.}
\begin{enumerate}[resume]
\item Are election contracts gambling or hedging instruments?
\item What did the 2024 court decision decide, and what did it leave open?
\item Which venue's price is the better forecast, and how would you know?
\item State the \emph{named result}: the price gap that survives fees and capital, and the Kelly stake on the cheaper contract for a belief of 65\%.
\item In one sentence: what is a price of 0.58?
\end{enumerate}
\end{problem}

\section{Interview questions}

\begin{interviewq}[$\star$ \iqroles{trader}]\label{iq:m3:prediction-and-betting-markets:1}
How do you turn a \omterm{def:m3:prediction-and-betting-markets:book}{bookmaker}'s odds into probabilities?
\end{interviewq}

\begin{interviewq}[$\star$ \iqroles{researcher}]\label{iq:m3:prediction-and-betting-markets:2}
What is the \omterm{def:m3:prediction-and-betting-markets:flb}{favourite--longshot bias} and what explains it?
\end{interviewq}

\begin{interviewq}[$\star\star$ \iqroles{trader}]\label{iq:m3:prediction-and-betting-markets:3}
How would you make a market in an \omterm{def:m3:prediction-and-betting-markets:event}{event contract} a month before the event?
\end{interviewq}

\begin{interviewq}[$\star\star$ \iqroles{risk}]\label{iq:m3:prediction-and-betting-markets:4}
What are the risks of holding offsetting positions in the same event on two venues?
\end{interviewq}

\begin{interviewq}[$\star\star$ \iqroles{researcher}]\label{iq:m3:prediction-and-betting-markets:5}
Derive the Kelly stake for a bet at decimal odds $o$ with win probability $p$.
\end{interviewq}

\begin{interviewq}[$\star\star\star$ \iqroles{developer, risk}]\label{iq:m3:prediction-and-betting-markets:6}
Design an \omterm{def:m3:blockchains-for-traders:contract}{oracle} for resolving \omterm{def:m3:prediction-and-betting-markets:event}{event contracts} that is hard to manipulate.
\end{interviewq}
